\DocumentMetadata{
  pdfversion=2.0,
  pdfstandard=ua-2,
  lang=en-US,
  tagging=on,
  tagging-setup={math/setup=mathml-SE,tabsorder=structure}
}
\documentclass[11pt,letterpaper]{article}
\usepackage[margin=0.68in,includefoot]{geometry}
\usepackage{newcomputermodern}
\usepackage{amsmath,amscd,mathtools}
\providecommand{\square}{\mdlgwhtsquare}
\usepackage[hidelinks]{hyperref}
\hypersetup{
  pdftitle={Analysis First-Year Exam, January 2016, Part 2},
  pdfauthor={University of Florida Department of Mathematics},
  pdfsubject={Graduate mathematics examination},
  pdfdisplaydoctitle=true,
  bookmarksopen=true
}
\setcounter{secnumdepth}{0}
\setlength{\parindent}{0pt}
\setlength{\parskip}{0.28em}
\setlength{\emergencystretch}{3em}
\setlength{\leftmargini}{2.15em}
\setlength{\leftmarginii}{2.65em}
\setlength{\leftmarginiii}{3em}
\renewcommand{\labelenumi}{\textbf{\theenumi.}}
\pagestyle{empty}
\raggedbottom
\allowdisplaybreaks
\newenvironment{examproblems}[1]{%
  \begin{enumerate}%
  \setcounter{enumi}{#1}%
  \setlength{\topsep}{0.35em}%
  \setlength{\partopsep}{0pt}%
  \setlength{\itemsep}{0.48em}%
  \setlength{\parsep}{0.18em}%
}{\end{enumerate}}
\newenvironment{examsubparts}{%
  \begin{enumerate}%
  \setlength{\topsep}{0.18em}%
  \setlength{\partopsep}{0pt}%
  \setlength{\itemsep}{0.16em}%
  \setlength{\parsep}{0.1em}%
}{\end{enumerate}}
\newenvironment{examinstructions}{%
  \par\begingroup\itshape
}{%
  \par\endgroup\vspace{0.18em}
}
\makeatletter
\renewcommand\section{\@startsection{section}{1}{0pt}%
  {-1.25ex plus -.25ex minus -.1ex}%
  {0.55ex plus .1ex}%
  {\normalfont\large\bfseries}}
\renewcommand{\@maketitle}{%
  \newpage\null\vskip -1.2em%
  {\footnotesize\bfseries
    \href{https://www.ufl.edu/}{University of Florida}, \href{https://math.ufl.edu/}{Department of Mathematics}\par}%
  \vskip 0.18em%
  \tagstructbegin{tag=Title}%
    {\LARGE\bfseries\href{https://gma.math.ufl.edu/exams/analysis-first-year/}{Analysis First-Year Exam}, January 2016, Part 2\par}%
  \tagstructend%
  \vskip 0.35em%
  \tagmcbegin{artifact}%
    \hrule height 1pt%
  \tagmcend%
  \vskip 0.55em%
}
\makeatother
\title{Analysis First-Year Exam, January 2016, Part 2}
\author{}
\date{}
\begin{document}
\maketitle
\thispagestyle{empty}
\begin{examinstructions}
Do exactly two problems from Part A and two problems from part B. Write solutions in a neat and logical fashion, giving complete reasons for all steps.
\end{examinstructions}
\section{Part A}
\begin{examproblems}{0}
\item
Let \(f_n:X\to\mathbb{R}\) be a uniformly convergent sequence of continuous functions on a compact metric space~\(X\). Prove that the set \(\{f_n\}\) is equicontinuous on~\(X\).
\item
Let \(\sum_{n=0}^{\infty} a_nx^n\) be a power series which converges for all \(x\in\mathbb{R}\). State and prove the theorem concerning the uniform convergence of the series on finite intervals \([a,b]\). Must the series converge uniformly on all of~\(\mathbb{R}\)? Prove, or give a counterexample.
\item
Suppose \(f\geq 0\), \(f\) is continuous on \([a,b]\), and \(\int_a^b f(x)\,dx=0\). Prove that \(f(x)=0\) for all \(x\in[a,b]\).
\end{examproblems}
\section{Part B}
\begin{examproblems}{0}
\item
Let \(f_n:[0,1]\to\mathbb{R}\) be a sequence of continuous functions. Prove that \(g=\limsup f_n\) and \(h=\liminf f_n\) are Lebesgue measurable.
\item
For each of the following, either prove or give a counterexample (\(m\) denotes Lebesgue measure on~\(\mathbb{R}\); assume all functions are measurable):
\begin{examsubparts}
\item[\textnormal{a)}]
If \(f_n\) is integrable for all~\(n\), \(f_n\to f\) uniformly on~\(\mathbb{R}\) and~\(f\) is integrable, then \(\int_{\mathbb{R}} f_n\,dm\to\int_{\mathbb{R}} f\,dm\).
\item[\textnormal{b)}]
If \(f_n\to f\) uniformly on \([0,1]\) and~\(f\) is integrable, then \(\int_0^1 f_n\,dm\to\int_0^1 f\,dm\).
\item[\textnormal{c)}]
Suppose \(f_n\geq 0\) and \(\int_0^1 f_n\,dm=1\) for all~\(n\). If \(f_n\to f\) pointwise, then \(\int_0^1 f\,dm\leq 1\).
\end{examsubparts}
\item
State the monotone convergence theorem and Fatou’s theorem, and use the former to prove the latter.
\end{examproblems}
\end{document}
